\begin{blocksection}
    %calculator implementation
    This code defines how a calculator represents expressions: an Expr is a variable, a number, or a Binary node that holds an operator and two smaller expressions, such that (1 + 2) * x becomes a tree. The helper functions check that a variable is a single letter, convert operators to symbols, and do arithmetic on two numbers.
    \begin{lstlisting}
    // The four arithmetic operators: +, -, *, and /.
    pub type Operator {
        Add
        Sub
        Mul
        Div
    }

    // An expression that can be evaluated.
    pub type Expr {
        Var(String)
        N(Int)
        Binary(operator: Operator, left: Expr, right: Expr)
    }

    // A variable name must be one letter from a to z.
    pub fn is_variable(name: String) -> Bool {
        string.length(name) == 1 && string.contains("abcdefghijklmnopqrstuvwxyz", name)
    }

    // The symbol for each operator.
    pub fn symbol(operator: Operator) -> String {
        case operator {
            Add -> "+"
            Sub -> "-"
            Mul -> "*"
            Div -> "/"
        }
    }

    // The result of applying an operator to numbers.
    pub fn calculate(operator: Operator, a: Int, b: Int) -> Int {
        case operator {
            Add -> a + b
            Sub -> a - b
            Mul -> a * b
            Div -> a / b
        }
    }
    \end{lstlisting}
\end{blocksection}